%!TEX root = ../hutbthesis_main.tex
\chapter{绪论}

\section{研究背景及意义}
随着工业 4.0、智能驾驶、康复医疗、虚拟现实和人机协同作业不断发展，人机交互技术也在向更自然、更贴近人体、更高精度和更柔顺的方向进步，传统的电机驱动、手柄操作、按键输入和视觉定位等交互方式，虽然在生产和生活中使用得比较普遍，但是在贴合人体、动作还原、操作柔和度和反应速度等方面都有明显不足，很难满足康复练习、精细操作、沉浸式驾驶模拟等较高要求场景里人机配合的需要，肌肉驱动的生物力学人机交互技术是以人体生理运动为核心，通过模仿肌肉收缩、关节配合和力量传递的过程，把人体想要做的动作直接变成外部设备的控制指令，它具有操作自然顺畅、反应快、适用性好、安全性高等特点，已经成为新一代人机交互研究里比较受关注的方向\cite{guntupalli2026comparing}。MuJoCo 是目前国际上常用的多刚体动力学仿真工具，可以高效完成肌肉骨骼建模、接触动力学计算和力反馈仿真，为生物力学交互系统的开发提供了稳定可靠的仿真环境。User-in-the-Box 开源框架又把生物力学人体模型、交互任务、视觉感知模块进行了统一打包，支持指向、跟踪、旋转操作等常见交互任务的快速搭建，明显降低了肌肉驱动交互系统的开发难度，在这样的背景下，把方向盘操控这种很有代表性的人机交互任务作为研究对象，依靠 User-in-the-Box 框架和 MuJoCo 引擎，设计并做出一套带有视觉引导功能的肌肉驱动生物力学人机交互系统，并对它的控制精准度、稳定性和易用性进行完整测试，是有比较重要的理论意义和实际意义的。本研究在理论层面具备一定的参考意义，相关研究成果能够完善肌肉发力、关节活动与机械运作之间的动态关联规律，也能改良肌肉骨骼模型与旋转操作设备的适配形式，丰富生物力学和人机交互交叉领域的基础理论内容；同时，本研究也能为MuJoCo工具与User-in-the-Box方法在交互仿真场景中的运用，提供一套可复用的参考方案。在实际应用层面，本次设计的系统拥有较多适用场景，可应用于智能驾驶模拟实训、上肢功能损伤人群的康复训练、VR沉浸式操控体验、工业协作机器人柔性控制等多个领域；该系统能够有效提升人机交互的自然度、操作精准度与整体作业的安全性，也可以降低设备的操作门槛与日常训练的成本。除此之外，这套系统还能够当作机器人工程、智能装备、生物力学等相关专业的教学实验设备，可用于开展课程实践教学与学生创新实训活动，能够助力人机交互领域的人才培育和技术落地推广，整体具备不错的应用范围与推广意义\cite{wang2026jiyu}。

\section{国内外研究现状}
\subsection{国内研究现状}
国内的很多研究都会用到 MuJoCo 这一类公开的仿真平台，高校和科研团队会结合具体场景做简单实用的设计，也会用成本不高的方式实现对应的功能，在人体肌肉骨骼模型完善、控制方式改良、康复训练和驾驶模拟等方面都做出了不少有用的成果，国内机器人、康复医疗、智能交互等行业发展速度较快，关注肌肉驱动和人机交互方向的团队也变得越来越多，相关的研究项目逐年变多，研究内容也更偏向实际使用的需要。不少团队不再只看重理论上的创新，他们更在意方法能不能用、好不好实现，也希望用简单省事、开销较小的路径，做出运行稳定、效果可以达到要求的肌肉驱动交互系统。哈工大、清华、北航等高校会把主要精力放在人体肌肉骨骼模型的制作和肢体运动模拟上，研究人员会一遍遍调整模型构造和运动相关的数值，让模拟出来的动作和真人肢体动作更加接近，他们还会优化关节活动范围、肌肉受力传递方式等内容，让抓取、旋转、推拉这类动作看起来更自然，也能给后续的交互控制环节提供更可靠的运行条件，工作人员还会不断调试仿真数值，让模型抖动、位置偏移、姿态异常等情况尽量减少，让整体运动过程更平稳、更接近真实情况\cite{zhang2026shenjing}。

上海交大、北理工、华南理工等高校会把研究重心放在控制算法的改良上自适应控制、模糊 PID、低成本信号处理等方面做出了有自身特点的成果，这些团队会针对肌肉信号容易变动、容易被外界影响、不同人之间差别较大等情况，提出多种改良后的算法，让控制的准确程度和反应速度得到提升，这些算法可以在配置一般的设备上运行，整体花费更低，可以使用的场景也更多，适合在普通实验环境里开展验证工作，不需要搭配价格较高的采集设备，也能帮助肌肉驱动交互技术被更多人使用。这些算法的改进目标大多是降低反应时间、缩小动作偏差、让运动过程更顺滑，在实验环境里可以得到比较理想的效果，也能给实际系统的开发提供可以参考的技术方式\cite{wang2026jiyu}。

哈工程、中山大学等单位会把肌肉驱动技术用在船舶驾驶模拟、康复机器人等真实场景中，技术可以模仿人体肌肉发力的状态来完成操作控制，驾驶训练的效果会得到提升，康复治疗的过程也会变得更安全、更舒适，肌肉驱动用在驾驶模拟中可以让操作感受更真实，也能提升练习者的沉浸体验，用在康复训练中可以温和带动肢体活动，能降低再次受伤的可能性，这些研究都给肌肉驱动交互系统用到实际场景打下了基础，也让相关技术慢慢从实验室走向可以真实使用的阶段，研究人员还能在真实场景测试中发现更多问题，这些问题在单纯仿真环境里不容易被察觉，可以据此继续完善模型和算法，让整个系统运行得更稳定。

但是从整体情况来看，国内相关研究还存在不少比较明显的问题，大部分研究只会针对某一个模型、某一种算法或者某一个场景展开，研究内容比较零散，整体的连贯性较差，各个模块之间缺少统一的对接方式，很多研究只会关注其中一个部分，要么只完善模型，要么只优化算法，很少把模型、控制、反馈、评测等内容整合到一起形成完整系统。同时，可以把肌肉驱动、视觉引导、方向盘精准控制结合起来的完整仿真系统并不多，相关研究数量少，成熟程度也不高，很多研究没有统一通用的框架作为支撑，可重复使用性和拓展能力较差，一个场景的研究成果很难快速用到其他场景中，整体研究效率也会受到影响。

除此之外，不少研究成果还停留在实验仿真阶段，和真实场景的结合程度不高，多种信息的融合效果也不够好，信号处理、视觉反馈、动作控制等环节的配合不够顺畅，数据不同步、反应不及时、反馈不清晰等情况经常出现，这些问题都会限制肌肉驱动人机交互系统在真实驾驶模拟、康复训练等场景中的使用和推广。国内研究在标准化、整体化、工程化等方面还存在明显不足，缺少像国外成熟框架那样完整的工具链条，研究人员需要花费大量时间处理基础问题，很难把主要精力放在关键技术创新上。总的来说，国内肌肉驱动与人机交互相关研究已经有不错的发展基础，也形成了自身的发展特点，但是在系统整合、场景适配、工程落地等方面还有较大的提升空间，这也给本文依托成熟开源框架进行方向盘操控和精度测试的研究，提供了实际的研究需要和研究意义。
\subsection{国外研究现状}
国外在生物力学建模、肌肉驱动仿真和人机交互领域的起步时间相对更早，整体的理论体系和技术方法都经过长期发展，变得更加成熟完善，大量基础工具和关键技术都由国外科研团队最先提出并不断完善。早期的研究人员会通过计算机仿真方式还原人体肌肉和骨骼的真实运动状态，相关研究会为后续肌肉驱动交互系统的设计提供稳定的基础，随着仿真技术和硬件设备的不断升级，国外相关研究也慢慢从简单动作模仿，转向更加真实、更加稳定的交互控制方向。MIT 团队推出的 MuJoCo 物理引擎，是该领域具有代表性的仿真工具，这款引擎可以高效处理多关节运动和接触碰撞问题，还能较好还原肌肉和关节之间的力传递过程，稳定的计算效果让它成为肌肉骨骼仿真研究中的常用工具，全球各地的研究团队都会借助这一引擎开展相关实验\cite{sitaram2026biomechanical}。在这类通用仿真平台的支持下，肌肉运动仿真的真实感和计算速度都得到明显提升，整个领域的研究推进速度也变得更快。

斯坦福大学、慕尼黑工业大学、苏黎世联邦理工学院等高校的科研团队，都会围绕人体肌肉骨骼模型搭建、肌肉信号处理、力反馈控制等方向开展大量研究工作，研究人员会不断优化模型结构，调整信号处理方式，改进对应的控制策略，交互过程的精准程度和动作的真实程度都会得到相应提升\cite{liu2026enhancing}。国外团队大多重视研究成果的通用性和可复用性，很多研究团队会把成熟的模型和算法打包成通用工具，这些共享的资源可以让更多研究者快速开展实验，也让整个领域的发展速度得到有效提升。经过持续的优化与改进，肌肉驱动控制系统的响应速度变得更快，控制误差可以进一步减小，运动过程的流畅性和稳定性也能逐步满足日常实验的基本要求。

国际团队搭建完成的 User-in-the-Box 框架，还可以把生物力学人体模型、交互任务、视觉感知等多个部分整合在一起，整套系统具备标准化的仿真流程和调用方式。该框架可以完成指向操作、轨迹跟踪、设备操控等多种常见的人机交互任务，还能结合强化学习方法自动优化交互策略，框架的整体结构、模块划分和运行逻辑，都能为本文的研究内容提供直接参考。这一框架降低了生物力学仿真的使用难度，研究者可以把更多精力放在控制逻辑和交互效果上，不用花费大量时间重复搭建基础模型，肌肉驱动交互相关的研究也可以更顺利地开展下去\cite{renganathan2026simulation}。

在方向盘、操纵杆等旋转类操控设备的交互研究中，国外团队已经推出将肌肉驱动、力反馈和视觉引导结合在一起的完整系统，这些系统可以完成稳定的角度跟踪和连续操控任务，仿真真实感和操作流畅度都能达到较好水平，还可以完成转向、角度保持、自动回正等一系列常规操作任务\cite{liu2026actuation}。但是这类系统也存在一定的短板与不足，系统整体的搭建成本和运行成本偏高，对硬件和软件的配置要求也相对更高，模型和操控设备之间的定制化配合程度不够理想，在更换不同实验场景时需要重新调整大量参数，无法做到快速切换与适配。部分系统过度追求高精度表现，实用性和普及性没有得到足够重视，普通实验环境很难完整复现相关实验，这些现存问题都为后续研究提供了可以改进和优化的方向。

从整体发展情况来看，国外相关研究起步早、积累充足、通用工具更加成熟，在模型精准程度、算法稳定程度和系统完整程度上都具备一定的领先优势，也为肌肉驱动人机交互领域提供了大量可以参考的思路与方案。同时，国外研究在平台搭建和标准制定方面的探索更加成熟，逐步形成完整的工具链和开发体系。但是系统成本偏高、定制化配合不足、场景适配不够灵活等问题依旧存在，这也为本文依托通用开源框架，搭建成本更低、运行更稳定、便于开展精度测试的方向盘肌肉驱动控制系统，提供了清晰的研究切入点和实际的研究意义。
\subsection{研究综述}

综合对国内外研究情况的整理可以发现，国外相关研究在基础模型搭建、算法运行效果、平台成熟程度上都有较好的表现，但是也存在通用性太高、针对性适配不够、实际部署花费较大等问题，国内研究则在实际场景使用、低成本实现、工程化应用等方面更有优势，不过在理论体系完整度、标准化水平和跨场景使用能力上还需要进一步提高，当前肌肉驱动人机交互系统大多存在不同人群适配效果不好、多种信息融合不够、操控精准度和稳定性不能同时保证等问题，面对这些研究中存在的不足，本课题以方向盘视觉引导操作为研究目标，借助 User-in-the-Box 框架和 MuJoCo 引擎，搭建一体化的肌肉驱动生物力学人机交互系统，重点改善系统的控制精准度、反应速度和使用便捷度，以此弥补现有技术存在的缺陷。
\section{研究内容与组织架构}
本研究把肌肉驱动的生物力学人机交互系统当作研究对象，主要围绕方向盘的视觉引导操作和精度测试进行探究，研究内容主要包含四个方面，一是借助 User-in-the-Box 框架和 MuJoCo 引擎，直接使用框架里自带的人体上肢肌肉骨骼模型与方向盘模型，完成模型的载入、位置调整和交互绑定，让肌肉激活信号可以带动方向盘进行转动，二是设计视觉引导模块和肌肉驱动控制算法，实现动态目标角度的生成、实时画面反馈，还有包含增益、死区、滤波、饱和限制在内的非线性对应关系，让操作过程更加顺畅和准确，三是运用 Python 语言完成仿真载入、肌肉控制、轨迹生成、数据收集、精度评测等模块的整合开发，搭建出可以稳定复现的完整实验系统，四是制定多方面的测试方案，开展功能测试、控制精度测试和稳定性测试，运用 MAE、RMSE、响应延迟、漂移量等指标对系统性能进行量化评测，并完成结果的整理与分析。

第 1 章是绪论，主要说明研究的背景与意义，整理国内外相关研究情况，确定研究内容和论文的整体结构。

第 2 章介绍相关的理论基础和关键技术，包括上肢生物力学、MuJoCo 仿真引擎、User-in-the-Box 框架、视觉引导以及精度评测指标等内容。

第 3 章进行系统的整体设计，开展需求分析，搭建模块化的系统架构，确定开发环境和使用工具。

第 4 章完成肌肉驱动生物力学模型的搭建与实现，包括人体肌肉骨骼模型、方向盘模型以及肌肉和方向盘耦合驱动的设计与检验。

第 5 章完成视觉引导与控制算法的设计和实现，包括算法开发、视觉界面制作和系统联合调试。

第 6 章进行系统测试与精度分析，制定实验方案并对测试结果进行对比和研究。

第 7 章是总结与展望，对研究成果进行归纳，指出研究的不足之处，并对今后的改进方向进行展望，如图~\ref{fig:structure}为本论文的结构框架图。
\begin{figure}[htbp]
    \centering
    \includegraphics[width=0.8\textwidth]{image_1.png}
    \caption{论文结构框架图}
    \label{fig:structure}
\end{figure}
